\chapter{Secure Development Practices}
\label{ch:secure-dev}

\epigraph{``Security is not a feature. Security is a property of the engineering process.''}{}

\learninggoal{
	\item Integrate security into the software development lifecycle for systems programming.
	\item Apply threat modelling to low-level systems (protocol parsers, device drivers, kernels).
	\item Use static and dynamic analysis tools effectively.
	\item Write and enforce secure coding standards for C and C++.
	\item Design test harnesses that catch memory-safety bugs before release.
}

\section{The Secure Development Lifecycle (SDL)}

For low-level software, security cannot be retrofitted. It must be embedded in each phase:

\begin{longtable}{p{2.5cm}p{5cm}p{5cm}}
	\toprule
	\textbf{Phase} & \textbf{Security activity}    & \textbf{Low-level focus}                        \\
	\midrule
	Requirements   & Threat modelling              & Identify trust boundaries, attacker model       \\
	Design         & Security architecture         & Privilege separation, minimal API surface       \\
	Implementation & Secure coding standards       & Bounds checking, safe API usage, format strings \\
	Testing        & Fuzzing + sanitizers          & Coverage-guided fuzzing with ASan/UBSan         \\
	Audit          & Code review + static analysis & Manual review of memory-unsafe patterns         \\
	Release        & Binary hardening              & Compiler flags, RELRO, PIE, seccomp             \\
	Response       & Incident plan                 & Patch process for memory-corruption CVEs        \\
	\bottomrule
\end{longtable}

\section{Threat Modelling for Low-Level Systems}

Threat modelling asks: what can an attacker control, and what can they achieve?

\subsection{Identifying Trust Boundaries}

For a system component, ask:

\begin{itemize}
	\item Where does the input come from? (network, file, user, another process, kernel)
	\item Is the input attacker-controlled? (yes for network parsers, maybe for file parsers, no for local configuration)
	\item What happens to the input? (parsed, copied, stored, executed)
	\item What privileges does this component have? (user, root, kernel)
	\item What happens if the component crashes? (DoS, privilege escalation, information leak)
\end{itemize}

\begin{definition}[Trust Boundary]
	A \textbf{trust boundary} is the point at which data crosses from a less-trusted to a more-trusted environment. Every crossing is a place where validation must occur.
\end{definition}

\subsection{The STRIDE Model}

Microsoft's STRIDE threat classification helps identify threat types:

\begin{longtable}{p{2.5cm}p{4cm}p{5cm}}
	\toprule
	\textbf{Threat}        & \textbf{Definition}           & \textbf{Low-level example}                 \\
	\midrule
	Spoofing               & Pretending to be someone else & Forged IP source address                   \\
	Tampering              & Modifying data or code        & Buffer overflow overwriting return address \\
	Repudiation            & Denying an action             & No audit trail memory corruption           \\
	Information disclosure & Leaking data                  & Format string leaking stack                \\
	Denial of service      & Making system unavailable     & Infinite loop in a protocol parser         \\
	Elevation of privilege & Getting more access           & Kernel exploit via ioctl                   \\
	\bottomrule
\end{longtable}

\section{Coding Standards for Memory Safety}

\subsection{Rule 1: Always Bounds-Check}

Every array access, every pointer arithmetic operation, every copy with a programmer-specified length must have a corresponding bounds check.

\begin{lstlisting}[language=C, caption=Bounds checking pattern]
int process_message(const uint8_t *buf, size_t buf_len) {
    if (buf_len < HEADER_SIZE)
        return -EINVAL;
    uint16_t payload_len = ntohs(*(uint16_t *)(buf + 2));
    if (payload_len > buf_len - HEADER_SIZE)  // CRITICAL: bounds check
        return -EINVAL;
    // Safe to copy payload_len bytes
    memcpy(payload_buf, buf + HEADER_SIZE, payload_len);
    return 0;
}
\end{lstlisting}

\subsection{Rule 2: Use Safe String Functions}

\begin{longtable}{p{3cm}p{5cm}p{4cm}}
	\toprule
	\textbf{Avoid}             & \textbf{Use instead}                & \textbf{Reason}                 \\
	\midrule
	\texttt{strcpy}            & \texttt{strlcpy}, \texttt{snprintf} & Always specify destination size \\
	\texttt{strcat}            & \texttt{strlcat}                    & Same                            \\
	\texttt{sprintf}           & \texttt{snprintf}                   & No buffer size argument         \\
	\texttt{gets}              & \texttt{fgets}                      & No bounds at all                \\
	\texttt{scanf("\%s", buf)} & \texttt{scanf("\%Ns", buf)}         & Specify max field width         \\
	\bottomrule
\end{longtable}

\subsection{Rule 3: Validate All Integer Arithmetic}

Integer overflows can mask buffer overflows:

\begin{lstlisting}[language=C, caption=Safe integer arithmetic]
#include <stdckdint.h>  // C23

bool safe_allocation(size_t count, size_t elem_size, void **out) {
    size_t total;
    if (ckd_mul(&total, count, elem_size))  // Checked multiplication
        return false;  // Overflow
    *out = malloc(total);
    return *out != NULL;
}
\end{lstlisting}

The \texttt{ckd\_*} functions (checked integer arithmetic) return true on overflow. Use them for every multiplication or addition that involves a length derived from external input.

\subsection{Rule 4: Avoid Variable-Length Arrays (VLAs)}

VLAs on the stack are dangerous because the size comes from a variable---if the variable is attacker-controlled, the stack may overflow:

\begin{lstlisting}[language=C, caption=Never use VLAs with untrusted sizes]
void process(size_t size) {
    // DANGEROUS: if size is large, this overflows the stack
    // and size is attacker-controlled
    char buf[size];
}
\end{lstlisting}

\subsection{Rule 5: Initialize All Memory}

Uninitialised memory may contain stale data from a previous allocation---including secrets or function pointers:

\begin{lstlisting}[language=C, caption=Zero-initialise allocated memory]
struct config *cfg = calloc(1, sizeof(*cfg));  // calloc zeros memory
// NOT: struct config *cfg = malloc(sizeof(*cfg)); // uninitialised!
\end{lstlisting}

\section{Fuzzing for Security Bugs}

Fuzzing is the single most effective technique for finding memory-safety bugs. Coverage-guided fuzzers (libFuzzer, AFL++) instrument the target binary and mutate inputs to maximise code coverage.

\begin{lstlisting}[language=C, caption=libFuzzer harness for a parser]
#include <stdint.h>
#include <stddef.h>

int LLVMFuzzerTestOneInput(const uint8_t *data, size_t size) {
    // Parse the input
    struct message msg;
    if (parse_message(data, size, &msg) < 0)
        return 0;  // Valid input rejected (not a bug)
    // Trigger deeper processing
    process_message(&msg);
    return 0;
}
\end{lstlisting}

Run with sanitizers:

\begin{lstlisting}[caption=Running the fuzzer with sanitizers]
$ clang -fsanitize=address,fuzzer -g -O1 fuzz_harness.c parser.c -o fuzz_parser
$ ./fuzz_parser -max_len=4096 -rss_limit_mb=2048 corpus/
\end{lstlisting}

Sanitizer integration is critical:
\begin{itemize}
	\item \textbf{AddressSanitizer:} detects heap/stack/global overflows, use-after-free, double free.
	\item \textbf{UndefinedBehaviorSanitizer:} detects integer overflow, misaligned access, null pointer dereference.
	\item \textbf{MemorySanitizer:} detects reads of uninitialised memory (requires all dependencies to be instrumented).
	\item \textbf{LeakSanitizer:} detects memory leaks.
\end{itemize}

\section{Static Analysis}

Static analysis tools find certain classes of bugs without running the code:

\begin{itemize}
	\item \textbf{Clang Static Analyzer:} path-sensitive analysis for null dereference, use-after-free, memory leaks.
	\item \textbf{Infer (from Facebook):} compositional analysis for large codebases.
	\item \textbf{CodeQL:} query-based analysis that can express complex security patterns (e.g., ``all format strings where the format argument is user-controlled'').
	\item \textbf{cppcheck:} simple pattern-based checker for common errors.
\end{itemize}

\begin{example}[Grep-Based Security Review]
	Before automated analysis, review for dangerous patterns:
	\begin{lstlisting}[caption=Grep patterns for security review]
# Find unbounded string copies
$ grep -rn 'strcpy\|strcat\|sprintf\|gets' src/

# Find unsafe format strings
$ grep -rn 'printf(user\|fprintf(user\|syslog(user' src/

# Find integer casts that might truncate
$ grep -rn '(int)\|(uint16_t)\|(uint8_t)' src/ | grep -v 'hton\|ntoh'
	\end{lstlisting}
\end{example}

\section{Compiler Hardening for Release Builds}

For production builds, enable all hardening flags:

\begin{lstlisting}[caption=Recommended compiler flags for security]
CFLAGS = -fstack-protector-strong \
         -fPIE -pie \
         -D_FORTIFY_SOURCE=3 \
         -Wl,-z,relro,-z,now \
         -Wl,-z,noexecstack \
         -fcf-protection=full \
         -mbranch-protection=standard
\end{lstlisting}

\section{Code Review Checklist for Low-Level Security}

A code review checklist for memory-safety-critical code:

\begin{enumerate}
	\item Every \texttt{memcpy}, \texttt{memmove}, \texttt{strcpy}, \texttt{snprintf} has a validated length argument.
	\item Every array index derived from external input is bounds-checked.
	\item Every integer arithmetic operation on lengths uses checked arithmetic (\texttt{ckd\_*} or manual overflow checks).
	\item Format strings are never user-controlled.
	\item All \texttt{switch} statements on external input have a \texttt{default} case.
	\item All structs copied from/to user space use \texttt{copy\_from\_user} / \texttt{copy\_to\_user} (kernel) or validate sizes (user).
	\item No variable-length arrays on the stack with attacker-controlled sizes.
	\item All memory containing secrets is explicitly zeroed with \texttt{volatile} or \texttt{explicit\_bzero}.
	\item Every allocation from \texttt{malloc} is checked for NULL.
	\item Function pointers are validated before calling (if feasible).
	\item All file descriptors are checked for validity before use.
	\item Every error path properly cleans up (frees memory, closes FDs, releases locks).
\end{enumerate}

\section{Building a Security Culture}

Technical measures are necessary but insufficient. The team must:

\begin{itemize}
	\item \textbf{Prioritise memory safety:} treat a memory-safety bug with the same severity as a production outage.
	\item \textbf{Invest in testing infrastructure:} continuous fuzzing with corpus management.
	\item \textbf{Conduct regular security reviews:} quarterly reviews of new code with the checklist above.
	\item \textbf{Adopt memory-safe components:} when feasible, use Rust or other safe languages for new protocol parsers, or isolate unsafe code behind safe APIs.
	\item \textbf{Track metrics:} number of memory-safety bugs found per release, time to fix, fuzzing coverage percentage.
\end{itemize}

\begin{principle}[The Ultimate Mitigation Is Process]
	Every mitigation---canaries, ASLR, CFI, seccomp---can be bypassed given enough skill and time. The only defence that scales is preventing exploitation from being necessary: find and fix bugs before they ship.
\end{principle}

\section{Exercises}

\begin{exercise}
	Apply the STRIDE threat model to a USB device driver in the Linux kernel. List at least one threat per category. What trust boundaries exist between user space and the driver?
\end{exercise}

\begin{exercise}
	Write a libFuzzer harness for a function that parses a simple binary protocol (e.g., a fixed header followed by a variable-length payload). Compile it with AddressSanitizer and run it for 1 million iterations. Report any crashes.
\end{exercise}

\begin{exercise}
	Audit the following code for security bugs. List every violation of the review checklist:
	\begin{lstlisting}[language=C]
int process_config(const char *input, size_t len) {
    char name[64];
    char *value = malloc(len);
    int n = sscanf(input, "%s=%s", name, value);
    if (n != 2) return -1;
    set_config(name, value);
    free(value);
    return 0;
}
	\end{lstlisting}
\end{exercise}

\begin{exercise}
	Research the Chromium sandbox architecture. How does it use privilege separation to contain memory-safety bugs in the rendering engine? What system calls are allowed in the sandboxed process?
\end{exercise}
